Guarded agentic compaction (\method) discovers and executes grounded initial
read\nobreakdash-only prefixes, with explicit admission assumptions and refusal
conditions, and leaves the rest of the agent unchanged. On three GitHub workflows the
compiled agent satisfies 90/90 held-out contracts against 89/90 for the
unchanged agent with 66.6\% fewer model requests, and a separately frozen
comparison with adapted offline Agent Workflow Memory over three workflows,
two models, and 1,362 fresh records (8,172 attempted episodes) cuts estimated
deployment cost by 36.9--79.1\% and requests by 50--75\%, with five of six
cells meeting a preregistered five-point accuracy noninferiority margin under
multiplicity correction. On BIRD, a schema-first SQL opening compiles on 26/31
databases with 45--50\% fewer requests and no detected accuracy loss, and a
post-cutoff cohort of 132 pull requests shows no end-to-end error (upper bound
1.73\% at 90\% confidence).

\method admits a program only when every tool argument is grounded in the entry
state, a prior result, or an invariant literal; every tool has a declared
read\nobreakdash-only effect; induced input/output contracts hold; and an exact
finite-sample test under independent-group assumptions bounds the dispatch
risk of the registered loss event. Otherwise it retires the candidate and keeps the agent, which is what
happens to question-specific table selection on BIRD, to every NESTFUL,
API-Bank, and BFCL family (too little calibration support), and, at dispatch
time, to AppWorld's admitted program on nearly all ReAct and plan-and-execute
runs. The evidence supports bounded compilation of recurrent prefixes, not
universal accuracy preservation or learned-gate benefit: most certificates
concern replay rather than final answers, the statistical stage adds no
measured runtime benefit on admitted candidates, one second-model cell is
inconclusive, and independence beyond the evaluated workloads is assumed.
